\chapter{LITERATURE REVIEW}%
\label{chap:literature_review}

Several studies and projects have explored the design and implementation of digital clocks and 
time-based embedded systems using microcontrollers. These works primarily focus on accurate time 
generation, efficient display control, and reliable system operation.

Microcontroller-based digital clocks are commonly implemented using hardware timers and 
interrupt mechanisms. According to Mazidi et al.~\cite{mazidi2014avr}, AVR microcontrollers 
provide flexible timer modules that can be configured in various modes, including Clear Timer 
on Compare (CTC) mode, which is suitable for precise time interval generation. The use of 
interrupts ensures that timing accuracy is maintained independently of the main program 
execution.

Seven-segment displays are widely used in embedded applications due to their simplicity, low 
cost, and ease of interfacing. Barret and Pack~\cite{barret2012avr} describe multiplexing as an 
efficient technique for controlling multiple seven-segment displays using a limited number of 
I/O pins. By rapidly switching between displays, multiplexing creates the illusion of continuous 
illumination while reducing hardware complexity.

Several embedded system designs utilize internal oscillators to reduce external component 
count and simplify circuit design. As discussed by the Atmel datasheet~\cite{atmega32datasheet}, 
the internal RC oscillator of the ATmega32 microcontroller is suitable for applications where 
ultra-high timing precision is not critical, such as basic digital clocks and timers.

Simulation tools such as Proteus have been extensively used for validating embedded system 
designs before hardware implementation. Proteus allows real-time simulation of 
microcontroller-based circuits, enabling developers to test both firmware and hardware 
interaction~\cite{proteus}. Similarly, Microchip Studio provides a robust development 
environment for AVR microcontrollers, supporting embedded C programming, debugging, 
and code optimization~\cite{microchipstudio}.

Based on the reviewed literature, it is evident that implementing a digital clock using AVR 
microcontrollers, hardware timers, interrupts, and multiplexed seven-segment displays is a 
well-established and reliable approach. This project follows proven design methodologies while 
serving as a practical learning platform for embedded system development.
